\documentclass[a4paper,11pt]{ltjsarticle}
\usepackage[haranoaji]{luatexja-preset}
\usepackage{amsmath,amssymb}
\usepackage[top=12mm,bottom=10mm,left=14mm,right=14mm]{geometry}
\usepackage{array}
\usepackage{tikz}
\usetikzlibrary{calc}
\pagestyle{empty}
\setlength{\parindent}{0pt}
\newlength{\qcol}\setlength{\qcol}{0.47\textwidth}
\newlength{\qgap}
\newlength{\anscolw}\setlength{\anscolw}{0.30\textwidth}
\newlength{\ansh}
\newcommand{\Q}[2]{\parbox[t]{\qcol}{\raggedright (#1)\hspace{0.6em}$\displaystyle #2$}}
\newcommand{\QR}[4]{\Q{#1}{#2}\hfill\Q{#3}{#4}\par\vspace{\qgap}}
\newcommand{\anscell}[1]{\rule[-\ansdrop]{0pt}{\ansh}(#1)}
\newlength{\ansdrop}

\begin{document}
{\large\bfseries 計算問題　29}\hfill 名前(\underline{\hspace{32mm}})\par\vspace{3mm}
■（1）〜（16）の計算をしなさい。（17）〜（21）は方程式を解きなさい。\par\vspace{3mm}
\setlength{\qgap}{5.4mm}
\QR{1}{(-14)+(-7)-(-13)}{2}{-7^2+(-4)^2-8}
\QR{3}{(7a+4)+(8a-2)}{4}{8\times\{4-(-8+1)\times2\}}
\QR{5}{\dfrac{7}{4}\div8+\dfrac{2}{4}}{6}{(-28)\div(+4)}
\QR{7}{(4x-7)-(2x+8)}{8}{\dfrac{x-7}{4}\times32}
\QR{9}{-0.7-(-8.4)-1.2}{10}{7\div4\times(-32)}
\QR{11}{\dfrac{7}{4}(4x-8)-\dfrac{8}{2}(2x+14)}{12}{15-(-4)\div\left(-\dfrac{4}{8}\right)}
\QR{13}{(28a-32)\div4}{14}{\dfrac{1}{7}(28x-56)+\dfrac{1}{2}(16x+8)}
\QR{15}{(-11)\times(-6)}{16}{7(4x-2)-(8x-2)}
\QR{17}{7x-4=17}{18}{7(x+4)=8x+23}
\QR{19}{0.7x+0.4=3.2}{20}{\dfrac{2x-3}{5}=\dfrac{x+2}{3}}
\vspace{1mm}
\Q{21}{7(x-2)-4(x+1)=11x-2}\par\vspace{3mm}
\setlength{\ansh}{9.5mm}\setlength{\ansdrop}{6mm}%
\begin{center}\begin{tabular}{|p{\anscolw}|p{\anscolw}|p{\anscolw}|}\hline
\anscell{1} & \anscell{2} & \anscell{3} \\\hline
\anscell{4} & \anscell{5} & \anscell{6} \\\hline
\anscell{7} & \anscell{8} & \anscell{9} \\\hline
\anscell{10} & \anscell{11} & \anscell{12} \\\hline
\anscell{13} & \anscell{14} & \anscell{15} \\\hline
\anscell{16} & \anscell{17} & \anscell{18} \\\hline
\anscell{19} & \anscell{20} & \anscell{21} \\\hline
\end{tabular}\end{center}

\end{document}
